\documentclass[11pt,a4paper]{article}
\usepackage[margin=1in]{geometry}
\usepackage{hyperref}
\usepackage{enumitem}
\usepackage{graphicx}
\usepackage{xcolor}

% -----------------------------------------------------
% bvraghav-tiet-ucs503p-article.sty
% -----------------------------------------------------

% Variable: Lab Instructor
\makeatletter \newcommand{\BvrSetLabInstructor}[1]{%
  \gdef\Bvr@LabInstructor{#1}}
% default
\newcommand{\Bvr@LabInstructor}{Dr. Raghav %
  B. Venkataramaiyer}
% public accessor
\newcommand{\BvrLabInstructorValue}{\Bvr@LabInstructor}
\makeatother

% Add subtitle
\usepackage{titling}
% -- Set title bold
\pretitle{\begin{center}\bfseries\fontsize{18bp}{18bp}\selectfont}
\makeatletter
\newcommand{\BvrSubtitle}[1]{%
  \posttitle{%
    \par\end{center}
    \begin{center}\large#1\end{center}
    \vskip 0.5em}%
}
\makeatother

% -----------------------------------------------------

\title{Project SocMag : Intelligent Society Recruitment \& Interview Portal}

\BvrSetLabInstructor{Dr. Raghav B. Venkataramaiyer}

\BvrSubtitle{%
  Project Proposal for UCS503P  \\
  Submitted to: \BvrLabInstructorValue \\ \bfseries
  Thapar Institute of Engineering and Technology%
}

\author{
  Member 1: Aishna Jindal%
  \thanks{Role: Frontend UI/UX \& Visualizations. Roll No: \texttt{1024170247}, %
    Email: \texttt{aaishna\_be24\@thapar.edu}}
  \and 
  Member 2: Muskaan Singhal%
  \thanks{Role: Database Architecture \& BCNF Optimization. Roll No: \texttt{1024170248}, %
    Email: \texttt{msinghal\_be24\@thapar.edu}}
  \and 
  Member 3: Shivani Shrivastava%
  \thanks{Role: Math Engine \& Backend Microservice. Roll No: \texttt{1024170278}, %
    Email: \texttt{sshrivastava\_be24\@thapar.edu}}
}

\date{\today}

\begin{document}
\maketitle

\section{Higher-order goal}
This project aims to optimize student extracurricular alignment and talent acquisition in higher education institutions by eliminating guesswork during campus recruitment drives. While student placement and organizational management are broad university domains, the immediate engineering objective is to translate \emph{data-driven applicant matching and queue theory} into a deployable software platform.

\section{Time-to-value}
\begin{itemize}[leftmargin=0.7cm]
    \item We will deliver meaningful increments quickly by prototyping the core fresher onboarding flow (``vibe check'') and score output within a short iteration window.
    \item We will prioritize fast feedback over perfection by using weekly agile sprints backed by CI automation to reduce release friction and validate recruiter interfaces early.
    \item Success is defined by deploying a functional matching and audition queue prototype in Sprint 1, then refining statistical precision based on initial user trials.
\end{itemize}

\section{Problem Statement}
During annual university recruitment drives, first-year students blindly apply to dozens of societies based on peer pressure rather than genuine alignment. Simultaneously, core committees for major student bodies (e.g., CCS, MUDRA) are overwhelmed with filtering 800+ raw applications, causing severe scheduling friction and administrative burnout. Common pain points include:
\begin{itemize}[leftmargin=0.7cm]
    \item \textbf{Unfiltered Applicant Floods:} Students apply randomly without knowing society culture or skill fit.
    \item \textbf{Recruiter Burnout:} Executive bodies manually process hundreds of unstructured resumes and forms.
    \item \textbf{Waiting Room Chaos:} Inefficient physical audition queues cause massive delay and student drop-offs.
    \item \textbf{Lack of Data Integrity:} Unstandardized feedback leads to subjective, biased selection across rounds.
\end{itemize}

\textbf{Why this matters now (contextual relevance):} In a dense institutional setting with tight recruitment schedules, automated matching and queue modeling significantly reduce candidate evaluation time and eliminate audition room congestion.

\section{Proposed Solution}
\subsection{Overview}
Build a web-based platform titled \textbf{Project SocMag}—an intelligent compatibility engine and audition management system featuring:
\begin{itemize}[leftmargin=0.7cm]
    \item \textbf{Interactive Onboarding: } Gamified onboarding flow converting student traits into feature vectors.
    \item \textbf{Statistical Match Engine:} Maximum Likelihood Estimation (MLE) engine comparing applicants against society profiles.
    \item \textbf{Recruiter Dashboard:} Swipeable Kanban board for multi-round audition management.
    \item \textbf{Dynamic Queueing:} Queue management using Poisson and Binomial models to fast-track candidates.
\end{itemize}

\subsection{Core workflow}
\begin{enumerate}[leftmargin=0.7cm]
    \item Applicant completes the interactive ``vibe check'' onboarding to generate normalized feature vectors.
    \item Statistical engine runs MLE to compute match scores, rendered visually via animated Spider/Radar charts.
    \item Recruiter reviews ranked applicants on a Kanban dashboard and manages round transitions.
    \item Poisson/Binomial queue algorithms model arrival rates and success probabilities to dynamic-schedule physical interviews.
\end{enumerate}

\subsection{Operational constraints for an educational campus setting}
\begin{itemize}[leftmargin=0.7cm]
    \item Design for dynamic, low-latency responsiveness during peak campus recruitment surges.
    \item Decouple mathematical execution into an isolated analytical microservice to preserve core web performance.
    \item Focus on deployability using standard containerization and cloud-managed databases.
\end{itemize}

\section{Solution Approach}
This is an engineering project emphasizing robust software architectural design, statistical modeling, and relational database normalization.

\subsection{Statistical \& Mathematical Engine}
Applicant matching and queue optimizations are backed by formal statistical methods:
\begin{itemize}[leftmargin=0.7cm]
    \item \textbf{Compatibility Vectoring:} Raw inputs (projects, personality traits, preferences) are transformed into normalized vector representations.
    \item \textbf{Maximum Likelihood Estimation (MLE):} Computes the statistical likelihood of an applicant fitting a society's ideal candidate profile.
    \item \textbf{Queue Theory (Poisson \& Binomial):} Uses Poisson distribution to model applicant physical arrivals and Binomial distribution to project multi-round pass rates for candidate fast-tracking.
\end{itemize}

\subsection{Web \& Data Architecture}
A decoupled multi-tier architecture:
\begin{itemize}[leftmargin=0.7cm]
    \item \textbf{Frontend:} Next.js (React) with Framer Motion transitions and Chart.js radar visualizations.
    \item \textbf{Backend API:} Node.js (Express) handling state management and REST workflows.
    \item \textbf{Analytics Microservice:} Isolated Python/R service executing SciPy/NumPy calculations.
    \item \textbf{Database:} PostgreSQL schema strictly normalized to BCNF (Students, Societies, Audition\_Rounds, Rubric\_Scores).
\end{itemize}

\subsection{SE Lifecycle, Integration, \& CI/CD}
Software engineering practices include:
\begin{itemize}[leftmargin=0.7cm]
    \item \textbf{Agile/XP Sprints:} Scrum velocity tracking with structured User Stories.
    \item \textbf{Formal Design artifacts:} Data Flow Diagrams (DFDs), Entity-Relationship (ER) diagrams, and BCNF schema validation.
    \item \textbf{Information Hiding \& Testing:} MVC architecture ensuring controller-math separation; white-box unit testing for statistical functions and automated CI pipelines.
\end{itemize}

\section{Evaluation Criterion}
We measure project success using quantitative metrics derived from recruitment operational data.

\subsection{Primary evaluation metric}
\textbf{Audition Queue Processing Time (AQPT):} Time taken from candidate queue check-in to interview completion.
\begin{itemize}[leftmargin=0.7cm]
    \item Target: Reduction of median waiting room time by $\ge 35\%$ during pilot testing.
    \item Attribution: Measured using microsecond database timestamps logging candidate arrival vs. interview exit.
\end{itemize}

\subsection{Secondary evaluation metrics}
\begin{itemize}[leftmargin=0.7cm]
    \item \textbf{Match Score Accuracy:} Correlation between MLE compatibility scores and high recruiter rubric marks.
    \item \textbf{Recruiter Triage Speed:} Reduction in screening time per application using the swipeable Kanban UI.
    \item \textbf{System Uptime \& Latency:} API sub-second response times under simulated peak recruitment traffic.
\end{itemize}

\subsection{Pilot validation plan}
\begin{itemize}[leftmargin=0.7cm]
    \item Deploy pilot across 50--100 first-year students and 2 campus societies (e.g., CCS and MUDRA).
    \item Compare recruitment throughput and user satisfaction ratings against historical baseline processes.
\end{itemize}

\section{Scalability}
The solution scales both theoretically and operationally to handle institute-wide recruitment drives.

\begin{enumerate}[leftmargin=0.7cm]
    \item \textbf{Scalable (theoretically):} Supports hundreds of concurrent applicants and recruiters by:
    \begin{itemize}[leftmargin=0.7cm]
        \item decoupling the calculation engine from web API handling,
        \item enforcing BCNF database indexing for rapid queries during peak hours,
        \item using stateless JWT-based RESTful API design.
    \end{itemize}
    
    \item \textbf{Operationally deployable (practically):} Runs efficiently on standard cloud platforms:
    \begin{itemize}[leftmargin=0.7cm]
        \item containerized microservices managed via Docker,
        \item cloud PostgreSQL database hosting,
        \item automated continuous deployment pipelines via GitHub Actions.
    \end{itemize}
\end{enumerate}

\section{Engine availability heuristic}
Mature development stacks exist to rapidly build and deploy this system:
\begin{itemize}[leftmargin=0.7cm]
    \item Next.js and Node.js (Express) for reactive frontend and API management.
    \item R / Python (SciPy, NumPy) runtime environment for MLE and probability modeling.
    \item PostgreSQL relational DBMS with BCNF-compliant schema execution.
    \item Chart.js and Framer Motion for interactive visual analytics.
    \item GitHub Actions for automated unit testing and staging builds.
\end{itemize}

\section{Project scope and deliverables}
\subsection{Initial deliverable}
\begin{itemize}[leftmargin=0.7cm]
    \item Onboarding flow collecting student traits and basic vector conversion.
    \item Initial PostgreSQL BCNF relational schema deployment.
    \item Recruiter candidate overview list with basic status transition controls.
    \item Automated CI pipeline enforcing unit testing on API endpoints.
\end{itemize}

\subsection{Subsequent deliverables}
\begin{itemize}[leftmargin=0.7cm]
    \item Isolated analytical microservice running MLE compatibility matching.
    \item Poisson/Binomial queueing algorithms driving audition fast-tracking.
    \item Interactive Spider/Radar chart visualizer for student profiles.
    \item Swipeable recruiter Kanban board and analytics dashboard.
\end{itemize}

\section{Risks and mitigations}
\begin{itemize}[leftmargin=0.7cm]
    \item \textbf{Mathematical Engine Latency:} Cache computed vectors and run heavy statistical operations asynchronously.
    \item \textbf{Bias in Ideal Profiles:} Allow society executive boards to dynamically calibrate attribute weights.
    \item \textbf{User Adoption Barrier:} Maintain intuitive, dark-mode visual design with zero onboarding friction.
\end{itemize}

\section{Summary}
Project SocMag presents an intelligent web platform designed to eliminate campus recruitment chaos. By combining Maximum Likelihood Estimation, Poisson/Binomial queue theory, BCNF database engineering, and modern web frameworks, the system replaces blind application floods with data-driven matching. Developed using rigorous XP agile methodologies, CI/CD automation, and clear performance metrics, SocMag delivers measurable efficiency gains for both applicants and executive committees.

\end{document}
